%% kapa: exam=1 round=25 year=2014 period=1 form=A subject=civil_law subject_ko=민법 seq=15 number=15
%% source: Q-Net 감정평가사 기출문제 (cst00309) artlSeq=5101207 file=extracted/1008/5101207/2014년 제25회 감정평가사 1차시험 기출문제/1교시(A형).hwp
%% conversion: kordoc 4.24.51 -> tools/gen1.py (auto); review=proofread-sonnet; flags=none
\begin{question}{15}
법률행위의 해석에 관한 설명으로 옳지 않은 것은? (다툼이 있으면 판례에 의함)
\begin{choices}
  \item 오표시무해(誤表示無害)의 원칙은 법률행위 해석 중 자연적 해석에 따른 것이다.
  \item 비전형의 혼합계약을 해석함에는 사용된 문언의 내용에 의하여 당사자가 그 표시행위에 부여한 객관적 의미를 있는 그대로 확정하는 것이 필요하다.
  \item 당사자가 특정 토지를 계약 목적물로 합의하였으나 그 지번의 표시에 관한 착오로 인하여 계약서에 그 토지와 다른 토지로 표시한 경우, 계약서에 표시된 토지에 대하여 계약이 성립한다.
  \item 사실인 관습은 당사자의 주장에 구애됨이 없이 법원이 스스로 직권에 의하여 판단할 수 있다.
  \item 계약을 체결한 자가 타인의 이름으로 법률행위를 한 경우, 계약당사자의 확정에 관한 행위자와 상대방의 의사가 일치하면 그 일치한 의사대로 행위자 또는 명의인을 계약의 당사자로 확정해야 한다.
\end{choices}
\end{question}
\answer{3}
